% !TEX root = ../../Extensions of IQC Duality to Partial Integral Equations.tex
In this work we assume that the generalized plant $P$ is a PIE system, that admits a state-space representation.
This work builds on the duality results in \cite{shivaknew} which shows that the state-feedback controller results in a convex optimization problem.
We find that this duality result can be extended to the case of dynamic filters, where we want to know the properties of the dual dynamic filter. This is particularly
usefull when examining the effect of uncertainties on the system when performing state-feedback controller synthesis, when using IQCs. The works from
Seiler and Venkataraman \cite{6915700,venkataraman} show that, after $J$-spectral factorization, the dual of a dynamic filter can be found using the Potapov-Ginzburg transform.
This section extends these results to PIE systems, and shows that the dual of a PIE system with a dynamic filter can be found
using the Potapov-Ginzburg transform, identical to the work from Venkataraman. 

\subsection{Duality under Dynamic Filtering}
We begin by restating the main theorem from \cite{shivaknew}.

\begin{figure*}[t]
\centering
\newcommand{\PrimalFilteredSymbols}{%
    \def\fdDelta{\Delta}%
    \def\fdPlant{G}%
    \def\fdFilter{\hat{\Psi}}%
    \def\fdZ{z}%
    \def\fdTZ{\tilde{z}}%
    \def\fdTW{\tilde{w}}%
    \def\fdW{w}%
}
\newcommand{\DualFilteredSymbols}{%
    \def\fdDelta{\Delta_d}%
    \def\fdPlant{G^\top}%
    \def\fdFilter{\mbf{D}(\hat{\Psi})}%
    \def\fdZ{\underline{z}}%
    \def\fdTZ{\tilde{\underline{z}}}%
    \def\fdTW{\tilde{\underline{w}}}%
    \def\fdW{\underline{w}}%
}
\newcommand{\PrimalPotapovSymbols}{%
    \def\pgFilter{\hat{\Psi}^{\dagger}}%
    \def\pgPlant{G}%
    \def\pgTW{\tilde{w}}%
    \def\pgTZ{\tilde{z}}%
    \def\pgZ{z}%
    \def\pgW{w}%
}
\newcommand{\DualPotapovSymbols}{%
    \def\pgFilter{(\hat{\Psi}^{\dagger})^\top}%
    \def\pgPlant{G^\top}%
    \def\pgTW{\tilde{\underline{w}}}%
    \def\pgTZ{\tilde{\underline{z}}}%
    \def\pgZ{\underline{z}}%
    \def\pgW{\underline{w}}%
}
\newcommand{\DrawFilteredPanel}[3]{%
\begin{scope}[shift={(#1,#2)}]
    \begin{scope}[shift={(3.05,1.52)}, scale=1.2, transform shape,
        shift={(-3.05,-1.52)}]
    \node[delta, dashed] (D) at (2.10,1.8) {$\fdDelta$};
    \node[block] (G) at (2.10,0.86) {$\fdPlant$};
    \node[block] (Psi) at (4,2.18) {$\fdFilter$};
    \coordinate (PsiInTop) at ([yshift=0.2cm]Psi.west);
    \coordinate (PsiInBot) at ([yshift=-0.2cm]Psi.west);
    \coordinate (PsiOutTop) at ([yshift=0.2cm]Psi.east);
    \coordinate (PsiOutBot) at ([yshift=-0.2cm]Psi.east);
    \coordinate (FdZLegBot) at ([xshift=-0.62cm]G.west);
    \coordinate (FdZLegTop) at (FdZLegBot |- D.west);
    \coordinate (FdWLegBot) at ([xshift=0.62cm]G.east);
    \coordinate (FdWLegTop) at (FdWLegBot |- D.east);
    \coordinate (FdZFilterTap) at ([xshift=0.3cm]FdZLegTop |- PsiInTop);
    \coordinate (FdWFilterTap) at ([xshift=-0.3cm]FdWLegTop |- PsiInBot);
    \draw[signal] (G.west) -- (FdZLegBot)
        -- node[siglabel, left] {$\fdZ$} (FdZLegTop) -- (D.west);
    \draw[signal] (D.east) -- (FdWLegTop)
        -- node[siglabel, right] {$\fdW$} (FdWLegBot) -- (G.east);
    \draw[signal] ([xshift=0.3cm]FdZLegTop) -- (FdZFilterTap)
        --  (PsiInTop);
    \draw[signal] ([xshift=-0.3cm]FdWLegTop) -- (FdWFilterTap)
        -- (PsiInBot);
    \draw[signal] (PsiOutTop) -- ([xshift=0.62cm]PsiOutTop)
        node[siglabel, right] {$\fdTZ$};
    \draw[signal] (PsiOutBot) -- ([xshift=0.62cm]PsiOutBot)
        node[siglabel, right] {$\fdTW$};
    \end{scope}
    \node[panelbox, minimum width=6.05cm, minimum height=3.08cm]
        (#3) at (3.05,1.52) {};
\end{scope}%
}
\newcommand{\DrawPotapovPanel}[3]{%
\begin{scope}[shift={(#1,#2)}]
    \begin{scope}[shift={(2.10,1.52)}, scale=1.2, transform shape,
        shift={(-2.10,-1.52)}]
    \node[block] (PsiPG) at (2.10,2.18) {$\pgFilter$};
    \node[block] (GPG) at (2.10,0.86) {$\pgPlant$};
    \coordinate (PsiPGInTop) at ([yshift=0.2cm]PsiPG.west);
    \coordinate (PsiPGInBot) at ([yshift=-0.2cm]PsiPG.west);
    \coordinate (PsiPGOutTop) at ([yshift=0.2cm]PsiPG.east);
    \coordinate (PsiPGOutBot) at ([yshift=-0.2cm]PsiPG.east);
    \coordinate (PgZLegBot) at ([xshift=-0.62cm]GPG.west);
    \coordinate (PgZLegTop) at (PgZLegBot |- PsiPGInBot);
    \coordinate (PgWLegBot) at ([xshift=0.62cm]GPG.east);
    \coordinate (PgWLegTop) at (PgWLegBot |- PsiPGOutBot);
    \draw[signal] ([xshift=-0.8cm]PsiPGInTop) node[siglabel, left] {$\pgTW$}
        -- (PsiPGInTop);
    \draw[signal] (GPG.west) -- (PgZLegBot)
        -- node[siglabel, left] {$\pgZ$} (PgZLegTop) -- (PsiPGInBot);
    \draw[signal] (PsiPGOutTop) -- ([xshift=0.8cm]PsiPGOutTop)
        node[siglabel, right] {$\pgTZ$};
    \draw[signal] (PsiPGOutBot) -- (PgWLegTop)
        -- node[siglabel, right] {$\pgW$} (PgWLegBot) -- (GPG.east);
    \end{scope}
    \node[panelbox, minimum width=4.48cm, minimum height=3.08cm]
        (#3) at (2.10,1.52) {};
\end{scope}%
}
\begin{tikzpicture}[
    >=Latex,
    block/.style={draw, minimum width=1cm, minimum height=1cm,
        align=center, inner sep=1pt, line width=0.5pt},
    delta/.style={draw, minimum width=0.6cm, minimum height=0.6cm,
        align=center, inner sep=1pt, line width=0.5pt},
    panelbox/.style={draw=none, fill=none, inner sep=0pt, outer sep=0pt},
    siglabel/.style={fill=white, inner sep=0.7pt, outer sep=1pt},
    signal/.style={->, line width=0.5pt},
    equiv/.style={
        {Implies[length=2.4mm]}-{Implies[length=2.4mm]},
        double,
        double distance=2pt,
        line width=0.5pt,
        shorten <=0.1cm,
        shorten >=0.1cm
    },
    every node/.style={font=\normalsize}
]
\PrimalFilteredSymbols
\DrawFilteredPanel{0.00}{4.15}{PrimalFilteredBox}
\PrimalPotapovSymbols
\DrawPotapovPanel{7.55}{4.15}{PrimalPotapovBox}
\DualFilteredSymbols
\DrawFilteredPanel{0.00}{0.00}{DualFilteredBox}
\DualPotapovSymbols
\DrawPotapovPanel{7.55}{0.00}{DualPotapovBox}

\draw[equiv] (PrimalFilteredBox.east) -- node[midway, below, yshift=-0.16cm]
    {(Lemma~\ref{lem:potapov-ginzburg-equivalence})} (PrimalPotapovBox.west);
\draw[equiv] (DualFilteredBox.east) -- node[midway, below, yshift=-0.16cm]
    {(Lemma~\ref{lem:potapov-ginzburg-equivalence})} (DualPotapovBox.west);
\draw[equiv] (PrimalFilteredBox.south) -- node[midway, right, xshift=0.16cm]
    {(Theorem~\ref{thm:filtered-dual-l2-gain})} (DualFilteredBox.north);
\draw[equiv] (PrimalPotapovBox.south) -- node[midway, right, xshift=0.16cm]
    {(Theorem~\ref{thm:dual-l2-gain})} (DualPotapovBox.north);
\end{tikzpicture}
\caption{Commutative diagram for primal and dual dynamic filtering under the
Potapov--Ginzburg transform.}
\label{fig:potapov-ginzburg-commutative-diagram}
\end{figure*}


\begin{theorem}[Dual $L_2$-gain
{\cite[Thm.~9]{shivaknew}}]\label{thm:dual-l2-gain}
For a PIE system with corresponding PI operators $P\{\mcl{T}, \mcl{A}, \mcl{B},
\mcl{C}, \mcl{D}\}$, and some $\gamma > 0$, the following statements are equivalent.
\begin{enumerate}
\item For any $w\in\mbf{L}_{\mathrm{e},[a,b]}^{\mbf{n}_w}$, and $\mbf{x}(0)=0$, $P$
satisfies $\|p_T {z}\|_{\mbf{L}} \leq \gamma \|p_T {w}\|_{\mbf{L}}$.
\item For any $\underline{w}\in\mbf{L}_{\mathrm{e},[a,b]}^{\mbf{n}_z}$, and
$\underline{\mbf{x}}(0)=0$, $P^\top$ satisfies $\|p_T \underline{z}\|_{\mbf{L}} \leq \gamma
\|p_T \underline{w}\|_{\mbf{L}}$.
\end{enumerate}
\end{theorem}
Relating an IQC to this gain duality requires two transformations. First,
write $\bmat{\tilde z\\\tilde w}:=\Theta\bmat{z\\w}$. A $J$-spectral
factorization gives, for every $T\geq0$,
\begin{equation*}
\begin{aligned}
&\ip{p_T\Psi\bmat{z\\w}}
{p_T\mcl M_{\mcl V}\Psi\bmat{z\\w}}_{\mbf L_{[a,b]}}\\
&\quad=
\|p_T\tilde z\|_{\mbf L_{[a,b]}}^2
-\|p_T\tilde w\|_{\mbf L_{[a,b]}}^2.
\end{aligned}
\end{equation*}
Thus the IQC quadratic form becomes the difference between the filtered
output and input signals. 

\begin{definition}[$J$-Spectral Factorization and Dual Filter
{\cite{venkataraman}}]
\label{def:dynamic-filters}\label{def:factorization}
Let $(\Psi,\mcl M_{\mcl V})$ be a hard IQC, let
\begin{equation*}
J_{\mbf n_z,\mbf n_w}:=
\bmat{I_{\mbf n_z}&0\\0&-I_{\mbf n_w}},
\end{equation*}
and let $\Theta,\Theta^{-1}, \Theta_{22}, \Theta_{22}^{-1}$ be causal bounded linear operators. The pair
$(\Theta,J_{\mbf n_z,\mbf n_w})$ is a $J$-spectral factorization of
$(\Psi,\mcl M_{\mcl V})$ if,
\begin{equation*}
\begin{aligned}
&
\ip{p_T\Psi\bmat{z\\w}}
{p_T\mcl M_{\mcl V}\Psi\bmat{z\\w}}_{\mbf L_{[a,b]}}
\\[-1mm]
&\quad=
\ip{p_T\Theta\bmat{z\\w}}
{p_TJ_{\mbf n_z,\mbf n_w}\Theta\bmat{z\\w}}_{\mbf L_{[a,b]}}.
\end{aligned}
\end{equation*}
Define the matrix
$\hat{J}:=\bmat{0_{\mbf n_z\times\mbf n_w}&I_{\mbf n_z}\\-I_{\mbf n_w}&0_{\mbf n_w\times\mbf n_z}}$. The dual filter associated with $\Theta$ is
defined by
\begin{equation*}
\mbf{D}(\Theta)
:=
\left(\hat{J}^*\Theta^\top\hat{J}\right)^{-1}.
\end{equation*}
\end{definition}
Give the primal inner product,
\[\mcl{I} := \ip{p_T \Theta\bmat{Pw \\ w}}
{p_T J_{\mbf{n}_z, \mbf{n}_w}\Theta\bmat{Pw \\ w}}_{\mbf L_{[a, b]}} \]
and corresponding dual iqc
\[\underline{\mcl{I}} := \ip{p_T \mbf{D}(\Theta)\bmat{P^\top\underline{w} \\ \underline{w}}}
{p_T J_{\mbf{n}_w, \mbf{n}_z}\mbf{D}(\Theta)\bmat{P^\top\underline{w} \\ \underline{w}}}_{\mbf L_{[a, b]}} \]
Although the following theorem can be derived using state-space realizations, the proof is cumbersome; for conciseness and clarity, we instead use an operator-based derivation.
\begin{theorem}\label{thm:dual-iqc}
    Let $P$ satisfy Definition~\ref{def:uncertain-interconnection} and let $(\Theta, J_{\mbf{n}_z,\mbf{n}_w})$ be a hard J-spectral factorization.
    Finally, assume $(\Theta_{21}P +\Theta_{22})$ has a bounded causal inverse. 
    Define 
    \begin{equation*}
        G_0 = (\Theta_{11}P +\Theta_{12})(\Theta_{21}P+\Theta_{22})^{-1}
    \end{equation*} 
    then the following statements are equivalent,
    \begin{enumerate}
        \item There exists a $0<\rho < 1$ such that $\|p_T G_0 \tilde{w}\|_{\mbf{L}_{[a,b]}}\leq \rho \|p_T \tilde{w}\|_{\mbf{L}_{[a,b]}}$ for every $\tilde{w}$ and $T\geq 0 $.
        \item There exists some $\epsilon > 0$ such that $\mcl{I}\leq - \epsilon \|p_T w\|_{\mbf{L}_{[a,b]}}^2$ for every $w$ and $T\geq0$.
        \item There exists some $\underline{\epsilon}>0$ such that $\underline{\mcl{I}} \leq - \underline{\epsilon}\|p_T\underline{w}\|_{\mbf{L}_{[a,b]}}^2$ for every $\underline w$ and $T\geq0$.
    \end{enumerate}
\end{theorem}
\begin{proof}
    See Appendix~\ref{proof:dual-iqc}.
\end{proof}

\begin{remark}[On the invertibility assumption]
Bounded causal invertibility of $\Gamma:=\Theta_{21}P+\Theta_{22}$ can be
certified using the bounded-real lemma. First verify that the feedthrough of
$\Gamma$ has a bounded causal inverse and use it to construct a PIE
realization of $\Gamma^{-1}$. Feasibility of the bounded-real KYP inequality
for this realization, with a coercive Lyapunov operator and a strict decay
margin, certifies exponential internal stability and finite $L_2$-gain of the
causal inverse.
\end{remark}

\begin{theorem}\label{thm:dual-KYP}
    Let the realization of $P$ satisfy Definition~\ref{def:uncertain-interconnection}. 
    Let $\Delta \in \mbf{\Delta}$ and there exist \textbf{either}, a primal hard IQC $(\Psi,\mcl{M}_{\mcl V})$,
    with $\mcl V\in\Pi_4$, or a dual hard IQC $(\underline{\Psi},\mcl{M}_{\underline{\mcl V}})$,
    with $\underline{\mcl V}\in\Pi_4$, such that,
    \begin{itemize}
        \item The interconnection $P\star\Delta$ is well-posed for all $\Delta\in\mbf{\Delta}$.
        \item The IQC in Definition~\ref{def:uncertainty-iqc} is satisfied for all $\Delta\in\mbf{\Delta}$.
        \item The primal hard IQC $(\Psi,\mcl{M}_{\mcl V})$ has a hard $(\Theta,J_{\mbf{n}_z,\mbf{n}_w})$ factorization.
    \end{itemize}
    Suppose that \textbf{either} of the two statements hold for some $\mcl{P} = \mcl{P}^* \succeq 0 \in \Pi_4$.
    \begin{itemize}
        \item \begin{minipage}[c]{\linewidth}
        \begin{equation*}
            \bmat{\mcl T^*\mcl P\mcl A + (*) & \mcl T^*\mcl P\mcl B\\
            \mcl B^*\mcl P\mcl T & \epsilon I} + \bmat{\mcl C^* \\ \mcl D^*}\mcl V\bmat{\mcl C & \mcl D} \preceq 0
        \end{equation*}
        \end{minipage}
        \item \begin{minipage}[c]{\linewidth}
        \begin{equation}\label{eq:dual-kyp}
            \bmat{\underline{\mcl T}^*\mcl P\underline{\mcl A} + (*)& \underline{\mcl T}^*\mcl P\underline{\mcl B}\\
            \underline{\mcl B}^*\mcl P\underline{\mcl T} & \underline{\epsilon} I} + \bmat{\underline{\mcl C}^* \\ \underline{\mcl D}^*}\underline{\mcl V}\bmat{\underline{\mcl C} & \underline{\mcl D}} \preceq 0
        \end{equation}
        \end{minipage}
    \end{itemize}
    Then for any $w,\underline{w}\in\mbf{L}_{\mathrm{e},[a,b]}$, satisfies either
    \begin{equation*}
        \Psi\bmat{P\\I}w := \left\{ \bmat{\mcl T \dot{\mbf{x}}(t) \\ v(t)} = \bmat{\mcl A & \mcl B\\
        \mcl C & \mcl D}\bmat{\mbf{x}(t) \\ w(t)}\right.
    \end{equation*}
    or
    \begin{equation}\label{eq:dual-system}
        \underline{\Psi}\bmat{P^\top\\I}\underline{w} := \left\{ \bmat{\underline{\mcl T} \dot{\underline{\mbf{x}}}(t) \\ {\underline{v}}(t)} = \bmat{\underline{\mcl A} & \underline{\mcl B}\\
        \underline{\mcl C} & \underline{\mcl D}}\bmat{\underline{\mbf{x}}(t) \\ \underline{w}(t)}\right.
    \end{equation}
    Then the interconnection $P\star\Delta$ with $\mbf{x}(0) = 0$ is robustly performing.
\end{theorem}
\begin{proof}
    See Appendix~\ref{proof:dual-KYP}
\end{proof}

\begin{remark}
    Note that we do not require information about the dual uncertainty $\underline{\Delta}$, only the primal uncertainty $\Delta$.
\end{remark}

